\documentclass[11pt]{article}
\usepackage[T1]{fontenc}
\usepackage[utf8]{inputenc}
\usepackage[margin=1in]{geometry}
\usepackage{amsmath,amssymb}
\usepackage[hidelinks]{hyperref}
\usepackage{url}
\setlength{\parindent}{0pt}
\setlength{\parskip}{0.55em}

\title{Response to the Associate Editor and Reviewers\\
\large Manuscript 26-0591v1: Informativity and Identifiability for Identification of Networks of Dynamical Systems}
\author{Anders Hansson, Jo\~ao Victor Galv\~ao da Mata, and Martin S.~Andersen}
\date{}

\begin{document}
\maketitle

We thank the associate editor and the three reviewers for their careful
assessment. The revised technical note distinguishes data informativity from
network identifiability, states the admissible model class more precisely,
provides proofs of the four theorems, and supplies a reproducible MATLAB
supplement. References to sections, theorems, and examples below refer to the
revised manuscript. The submitted manuscript's references [3], [17], and [19]
are, respectively, the papers by Hendrickx, Gevers, and Bazanella; van der
Woude; and Legat and Hendrickx. We address the comments individually below.

\section*{Response to the Associate Editor}

\textbf{Comment.}
\begin{quote}\itshape\emergencystretch=2em
Three reviewers evaluated the manuscript, and they converge on a set of substantive concerns regarding novelty, rigor, and presentation, which are summarized below:
\end{quote}

\textbf{Response.} We recognize the editor's assessment and the need for a major revision before resubmission. The revised note narrows its contribution claims, states its assumptions, and gives proofs and worked calculations for the results retained below.

\subsection*{AE.1: Novelty and positioning}
\textbf{Comment.}
\begin{quote}\itshape\emergencystretch=2em
1. Novelty and positioning of the results --- Reviewer 5; Reviewer 4. The main results appear to be straightforward extensions of well-known results. In particular, Reviewer 5 notes that the link between vertex-disjoint paths and the rank of associated matrices (Theorem 2) is already established in reference [3] of the manuscript. Reviewer 4 similarly observes that Theorem 2 "seems like a correct generalization" of a result in [17] but that the extension is not shown to be nontrivial, and asks whether Theorem 3 is new. The paper does not clearly delineate what is original versus what is a restatement of existing work, which must be clarified.
\end{quote}

\textbf{Response.} The revised abstract identifies the broader model class and the two-solution example. Three bullets at the end of Section I now state explicitly what is novel relative to the cited work: the scope of the joint analysis, its model-specific rank tests, and the identifiability analysis for the broader model with repeated and constrained entries. They identify the spectral, path-rank, and fiber-dimension principles as established results. Section IV applies the path result to $\Pi=RP^{-1}Q+S$ under independent free entries and gives a separate Gr\"obner basis rank test for known or constrained entries. Section V applies the finite-fiber criterion to the closed-loop map $F(X)=R_y(X)P(X)^{-1}Q_r(X)+S_{yr}(X)$, allowing unknown input and measurement transfer functions and repeated entries. Example 3 illustrates the analysis by computing two solutions that a generic Jacobian-rank test does not count.

\subsection*{AE.2: Proofs and references}
\textbf{Comment.}
\begin{quote}\itshape\emergencystretch=2em
2. Missing proofs or references for the main theorems --- Reviewer 3 ; Reviewer 4. Theorems 1--3 are given without formal proofs or precise references. Reviewer 3 notes that the informativity results are presented only informally as generalizations of prior work, and recommends inserting formal proofs stating exactly which assumptions are necessary. Reviewer 4 points out that Theorem 1 is loosely stated (the spectrum should be positive definite "for almost all \textbackslash{}omega"), and that references on the "transfer of excitation" problem (e.g., Gevers, Bazanella, Bombois, Miskovic, IEEE TAC, Vol 54, No 12) should be provided for context.
\end{quote}

\textbf{Response.} The revised manuscript includes a proof of Theorem 1 using the external-signal spectrum and the rank of $\Pi$, a determinant and path proof for Theorem 2, an elimination proof for Theorem 3 with the nonzero-denominator restriction, and a fiber-dimension proof for Theorem 4. The paper cites the underlying informativity, path-rank, and algebraic-geometry results.
Theorem 1 states the network, well-posedness, spectral, predictor-stability,
excitation, and generic-rank assumptions explicitly.

\subsection*{AE.3: Concepts and terminology}
\textbf{Comment.}
\begin{quote}\itshape\emergencystretch=2em
3. Conceptual clarity and terminology --- Reviewer 4; Reviewer 5; Reviewer 3. The manuscript mixes the notions of identifiability, network identifiability, generic/local/global identifiability without clear distinction. Reviewer 4 notes that the definition offered in the introduction conflates classical identifiability with network identifiability, and that the claim that partial measurement was not previously considered is incorrect. Reviewer 5 stresses that informativity and identifiability are distinct concepts that deserve separate treatment, and flags that Example 3 appears to contradict a conjecture discussed in [19], which should be carefully verified.
\end{quote}

\textbf{Response.} Sections I and V distinguish predictor informativity, generic network identifiability, and generic local network identifiability in the algebraic sense. The introduction now says explicitly that the latter permits several compatible networks, and Section V explains that the criterion does not certify an exceptional realization. Example 3 is no longer presented as a counterexample to the Legat--Hendrickx question: it uses unknown input and measurement transfer functions and a repeated unknown entry, which lie outside that paper's model class.

\subsection*{AE.4: Model and assumptions}
\textbf{Comment.}
\begin{quote}\itshape\emergencystretch=2em
4. Model formulation and assumptions --- Reviewer 5; Reviewer 3; Reviewer 5. The motivation for the model P(q)wk =Q(q)uk , zk =R(q)wk +S(q)uk is unclear: Reviewer 5 asks why two matrices P(q) and Q(q) are needed and notes the unusual separation of state and output dynamics compared to standard network models. Reviewer 3 points out that invertibility of P is assumed without stating conditions or references, that it is not stated whether the network topology is known, and that it is unclear how to select which predictors to use. Reviewer 5 further notes that the function F(X)=Ry P-1Qr +Syr is written as a function of X, yet X does not appear in the expression.
\end{quote}

\textbf{Response.} Section II embeds the familiar $w=Gw+Bu$, $z=Cw$ model as $P=I-G$, $Q=B$, $R=C$, $S=0$. It now motivates the separate $P$ and $Q$ through physical networks with local node dynamics, coupling laws, and a separate input matrix, citing Kivits and Van den Hof. It also gives a discrete-time descriptor model with $P(q)=qE-A$ as an illustration of a coefficient-constrained structure. Normalizing the node equations recovers an $I-G$ form but introduces shared unknown factors. Filtered or combined measurements and direct input effects motivate $R$ and $S$. Section II also states the known structure, separates generic invertibility from well-posedness, and checks $\det P$ in the examples. Section III explains how measured signals define the predictor.

\subsection*{AE.5: Presentation and readability}
\textbf{Comment.}
\begin{quote}\itshape\emergencystretch=2em
5. Presentation and readability --- Reviewer 5; Reviewer 4. Several sections (Section VI on sub-networks, and the numerical examples in Sections VII--VIII) are perceived as adding little and could be removed or substantially condensed. The paper relies on concepts from algebraic geometry (affine varieties, ideals, fibers, Grobner bases) that are not standard in the controls community and should be clearly introduced. The symbol z is used for two different quantities, and there are numerous typos (e.g., publication year "2913" in reference [21]).
\end{quote}

\textbf{Response.} We condensed the active subnetwork discussion, revised the examples around distinct uses of the rank and fiber criteria, and added short explanations of the algebraic terminology at its point of use. Section VIII explains when path, Jacobian-rank, and elimination tests are useful and notes the cost of symbolic elimination. The computational examples give software and hardware details and point to \path{supplementary_matlab_code.zip}. The shift operator is $q$, the complex argument is $\lambda$, and the output signal is $z_k$.

\textbf{Comment.}
\begin{quote}\itshape\emergencystretch=2em
Based on the collective reviews and this editor\textquoteright{}s own reading, this editor finds that the manuscript addresses a relevant topic and the reviewers acknowledge potential value in some ideas. However, the manuscript in its current form lacks formal justification for its main theorems, does not adequately position its novelty relative to the established literature, and contains unresolved conceptual and modeling ambiguities. These concerns go beyond minor revision, as the reviewers question whether the core results are nontrivial extensions. Based on these grounds, this editor suggest to reject this manuscript but with the option that it could be resubmitted as a technical note after major revision. If the authors decide to resubmit, the authors need to address all concerns raised by the reviewers in the revision, and moreover, the authors are reminded that IEEE TAC has the policy that the manuscript will be rejected if it cannot be accepted with minor revision in the next round of review.
\end{quote}

\textbf{Response.} We have revised the manuscript in response to the editor and all three reviewers. The detailed changes are described in the replies below.

\section*{Response to Reviewer 3}

\textbf{Comment.}
\begin{quote}\itshape\emergencystretch=2em
This manuscript deals with the informativity and identifiability conditions for network of dynamic systems. The authors focus their analysis using the so-called Grobner basis and provide graph-theoretic conditions. For informatitvity, sufficient conditions are derived in terms of the rank of a matrix of rational functions, relating the excitation signals to the inputs of the optimal predictor. The identifiability is investigated under the lens of the dimensionality of a proper fiber for the local identifiability. The paper assumes full excitation of the network, but partial measurement. This paper also provides an example under which it states that local identifiability does not imply generic identifiability, a hypothesis that was previously discussed in the literature. This paper is clear and well-written and there is also a clear contribution to the field, but it is hard to follow specially because the authors tend to omit key steps. Therefore, I have some concerns about the results of the paper given in the following.
\end{quote}

\textbf{Response.} We thank the reviewer for identifying where assumptions and steps needed to be stated explicitly.

\subsection*{R3.1: Invertibility and well-posedness of $P$}
\textbf{Comment.}
\begin{quote}\itshape\emergencystretch=2em
1. The network model used in this manuscript assumes that P is invertible for well-posedness. No remarks nor references are given to state under what conditions this transfer matrix is invertible.
\end{quote}

\textbf{Response.} Section II now assumes $\det P$ is not identically zero
as a rational expression in the free entries. This is a case-by-case test
for generic invertibility. A particular realization is called well-posed
only when its rational matrix $P$ is invertible; then $P(\lambda)$ is
invertible except at finitely many complex arguments where it is defined.
The examples display the relevant determinants. Stability is imposed
separately where needed.

\subsection*{R3.2: Formal proofs and necessary assumptions}
\textbf{Comment.}
\begin{quote}\itshape\emergencystretch=2em
2. There are no formal proofs for none of the Theorems presented in the informativity section. The ideas behind the proofs are simple given in informal terms, and mostly they are cited as generalizations of works from other. I recommend inserting formal proofs for all the results stating exactly which are necessary, instead of just referring to a book covering a wide topic.
\end{quote}

\textbf{Response.} Theorems 1--3 in Section IV now have proofs. Theorem 1
requires a positive-definite external-signal spectrum for almost all
frequencies and full generic row rank of $\Pi$; it concludes informativity
for algebraically generic admissible true networks, not for every realization.
It also states well-posedness of the true network, existence of the
external-signal spectrum, and stability of the admissible candidate predictors.
Theorem 2 permits prescribed nonzero diagonal entries of $P$ and requires
the other structurally nonzero entries to be independent free variables. It uses a
Schur-complement determinant identity to relate the rank of the transfer function matrix
to vertex-disjoint paths. Theorem 3 defines the denominator restriction
$D\ne0$, the minor ideal and the auxiliary equation $1-tD=0$, then proves
the three rank cases by elimination. A remark after Theorem 2 also covers
fixed selector matrices $Q,R$, formed from distinct columns and rows of
the identity matrix, respectively, with independent free structurally
nonzero entries in $S$.
More general known and constrained entries are handled by the algebraic rank test.
Example 1 contains fixed off-diagonal entries of $P$, so its informativity conclusion follows
from the displayed determinant of $\Pi$, rather than an application of
Theorem 2.

\subsection*{R3.3: Choice of predictor}
\textbf{Comment.}
\begin{quote}\itshape\emergencystretch=2em
3. While the characterization is given in terms of input to predictors input, how exactly one knows which predictors to use? There are several ways to select the predictors.
\end{quote}

\textbf{Response.} Section III makes the choice of measured output $y_k$
and measured input $r_k$ explicit through $z_k=(y_k,r_k)$. Under the assumed
stable innovation representation, this choice defines the predictor and the
closed-loop transfer function $G_c$ that it determines. Different choices
lead to different $G_c$ and require their own informativity test. The
predictor-reduction remark in Section IV and Example 2 show how known
relations among predictor columns can permit a smaller measured-input
vector. The input matrix $Q$ need not excite every node. We do not claim
a universal optimal predictor-selection rule.

\subsection*{R3.4: Known topology}
\textbf{Comment.}
\begin{quote}\itshape\emergencystretch=2em
4. Does one need to assume that the topology of the network is known? This assumption is not stated in the manuscript.
\end{quote}

\textbf{Response.} Section II now states that the structural zero pattern,
known entries, and specified relations among unknown entries of
$(P,Q,R,S)$ are given. The analysis determines unknown entries
within this model class; topology recovery is outside its scope. Theorem 2
requires independent free entries except for prescribed nonzero diagonal
entries of $P$; its following remark also permits fixed selector matrices
$Q,R$. The Gr\"obner basis procedure allows more general prescribed entries
and specified relations.

\subsection*{R3.5: Constructing the rank-test Gr\"obner basis}
\textbf{Comment.}
\begin{quote}\itshape\emergencystretch=2em
5. Theorem 3 seems interesting by its own, but it is not clear how to construct the Gronber basis from it.
\end{quote}

\textbf{Response.} Section IV gives the construction: clear row
denominators to obtain $\widetilde\Pi$, adjoin all $k$-minors and $1-tD$
to the ideal, place $t$ first in a lexicographic order, and inspect basis
polynomials independent of $t$ on $D\ne0$. Example 3 works through the
full-row-rank case: its reduced basis is $\{t-1,H\}$, so rank drops exactly
when $H=0$. The auxiliary variable excludes singular or otherwise
undefined realizations.

\subsection*{R3.6: Constructing the identifiability Gr\"obner basis}
\textbf{Comment.}
\begin{quote}\itshape\emergencystretch=2em
6. The same comment applies to Theorem 4, which derives a necessary and sufficient condition for local identifiability.
\end{quote}

\textbf{Response.} Section V describes how to clear denominators in
$F(X)=G_c$, adjoin $1-tD_F$, eliminate $X$ and $t$, and compute the
dimension of the resulting image closure. Example 3 shows the polynomial
equations, a lexicographic variable order, the zero elimination ideal,
and the resulting image dimension. It also solves the fiber equation to
show why equal dimensions can coexist with two admissible networks.

\section*{Response to Reviewer 4}

\textbf{Comment.}
\begin{quote}\itshape\emergencystretch=2em
The paper treats the identification of networks, particularly the problem of their identifiability, which is an important topic and has been receiving particular attention from the community in the past 15 years or so. It brings new and relevant ideas to the field and the results seems correct. Yet, there are various issues to be resolved before publication in these Transactions can be considered.
\end{quote}

\textbf{Response.} We thank the reviewer for the detailed comments on terminology, prior work, and proofs.

\subsection*{R4.1: Terminology for identifiability}
\textbf{Comment.}
\begin{quote}\itshape\emergencystretch=2em
In the introduction the concepts of informativity and identifiability are briefly presented. These are rather confused, as ``identifiability'' is a classical concept in the identification literature that usually refers to parametrizations of transfer functions, whereas the definition offered in the introduction (``whether the network topology is such that the open loop transfer functions can be uniquely determined from the predictor.'') looks more like what is called ``network identifiability'' in the literature of network identification. In the introduction itself the authors ``correct'' themselves in this regard when, in the last paragraph, they talk about local network identifiability instead. Please clarify.
\end{quote}

\textbf{Response.} The introduction now separates identification of a
predictor from recovery of the unknown open-loop transfer functions. The
latter is called \emph{network identifiability}, distinct from
identifiability of a numerical parameterization. Section V defines generic
local network identifiability as a zero-dimensional generic algebraic
fiber, which allows finitely many solutions; generic network
identifiability requires uniqueness. This finite-ambiguity meaning now
appears in the introduction as well. The manuscript also distinguishes
this symbolic rational-function formulation from the single-frequency
local notion in Legat and Hendrickx.

\subsection*{R4.2: Earlier work on partial measurements}
\textbf{Comment.}
\begin{quote}\itshape\emergencystretch=2em
It is also said that more general networks are considered than in previous works, which is correct. But It is also said ``we will consider the partial measurement case \ldots{}'' implying that previous works all considered the full measurement case, which is far from the truth. Please correct.
\end{quote}

\textbf{Response.} The introduction now cites prior work on partial node
measurements and subnetworks. We do not present partial measurement alone
as new. The model permits unknown input and measurement transfer
functions, specified relations among unknown entries, and a direct
input-to-output term; these are the aspects compared with prior work in
Sections I, II, and V.

\subsection*{R4.3: Standard MIMO informativity and Definition 1}
\textbf{Comment.}
\begin{quote}\itshape\emergencystretch=2em
Regarding informativity, what is studied is if the input signals are rich enough to identify the ``closed-loop'' relationship from u to z. So, this is standard MIMO identification, there is nothing ``networkly'' here, is it? Then definition 1 seems the standard definition in the literature, so some reference to it would be welcome; if this is not the standard definition, references are still needed, and the difference(s) should be commented. By the way, this is also the case in other parts of the paper, where some known results or definitions are stated without references; this too must be improved.
\end{quote}

\textbf{Response.} We agree that the positive-definite predictor-input
spectrum is a standard informativity condition, for which we now cite
Ljung's \emph{System Identification: Theory for the User}. Definition 1 is
attributed to the prediction-error informativity formulation of Bombois
et al.; the discussion also cites Gevers, Bazanella, Bombois, and
Miskovic for network-specific excitation. The manuscript explains why predictor uniqueness is the
relevant data condition before asking whether open-loop transfer functions are uniquely
recovered. The network-specific step is computing the generic rank of
$\Pi=RP^{-1}Q+S$ from the known interconnection structure or from
specified algebraic constraints. Informativity and network
identifiability are developed in separate sections.

\subsection*{R4.4: Distinct uses of $z$}
\textbf{Comment.}
\begin{quote}\itshape\emergencystretch=2em
One should not use \textquoteleft{}z\textquoteright{} to denote two different things.
\end{quote}

\textbf{Response.} The measured output signal retains the conventional
notation $z_k$. Transfer functions use $q$ as the shift operator and
$\lambda$ as a complex argument. The revised manuscript no longer uses
$z$ as a second transform variable.

\subsection*{R4.5: Theorem 1 and transfer of excitation}
\textbf{Comment.}
\begin{quote}\itshape\emergencystretch=2em
Now, the main issue: the Theorems. Theorems 1 to 3 are given without proof or reference! In Theorem 1 the statement itself is loose: it should probably read ``the spectrum is positive definite for almost all omega \ldots{}''. With this statement the Theorem is probably true, but it remains to be proven. Also, references on this ``transfer of excitation'' problem should be given for context and comparison, so that we (the readers) can tell whether the result is new; for that, I suggest that the authors look at the following paper from these Transactions and some of the references therein: Gevers, Bazanella, Bombois, Miskovic - Identification and the information matrix: how to get just sufficiently rich? Vol 54, No 12.
\end{quote}

\textbf{Response.} Theorem 1 now explicitly requires the external-signal
spectrum to be positive definite for almost all frequencies, meaning
outside a set of Lebesgue measure zero. Its proof uses
$\Phi_z=\Pi\Phi_u\Pi^*$ and full row rank of $\Pi$ for algebraically
generic admissible true networks. The citation to Gevers et al. provides
the requested transfer-of-excitation context, while Ljung is cited for
the classical spectral informativity implication. The manuscript presents
that implication as standard rather than a new network theorem.
The theorem also specifies well-posedness, existence of the
external-signal spectrum, and stability of the candidate predictors.

\subsection*{R4.6: Theorem 2 and the earlier path result}
\textbf{Comment.}
\begin{quote}\itshape\emergencystretch=2em
Theorem 2 indeed seems like a correct generalization of the very nice result in [17], but the extension is not trivial. Also, more context is welcome here, as similar disjoint paths conditions have been extensively used in the field. 
\end{quote}

\textbf{Response.} We cite van der Woude's structured-system result.
Theorem 2 allows prescribed nonzero diagonal entries of $P$, while all
other structurally nonzero entries must be independent free variables.
The revised proof uses
$\det\!\begin{bmatrix}P&Q_J\\-R_I&S_{I,J}\end{bmatrix}
=\det(P)\det(\Pi_{I,J})$. Determinant terms give disjoint input-output
paths, including one-edge paths through $S$. It now explains how the
permutation terms split into paths, cycles, and diagonal factors.
The entries selected outside the diagonal of $P$ determine the remaining
diagonal factors, so prescribed nonzero diagonal entries do not introduce
cancellation. A remark defines selector matrices $Q,R$ and explains why
the determinant argument also applies to them. Here $S$ may be zero or
have independent free entries. For $P=I-G$ and $S=0$, this recovers the
path-rank result for selected submatrices of $(I-G)^{-1}$ in Hendrickx,
Gevers, and Bazanella. We present the
general path/rank principle as prior work and specify how the present
matrix form is covered.

\subsection*{R4.7: Theorem 3, references, and novelty}
\textbf{Comment.}
\begin{quote}\itshape\emergencystretch=2em
Theorem 3 also needs a proof and/or references to provide better context. Is this Theorem new? Overall, the proofs of the Theorems - or references to the proofs if they are not original - must be given and more reference in the text to the papers in the reference list (which may be amended a little) must be used to provide better context.
\end{quote}

\textbf{Response.} Theorem 3 now has an elimination proof and an explicit
ideal construction; Example 3 gives a small calculation. We use
standard Gr\"obner basis elimination and cite the algebraic-geometry
references for that method. The claim is the rank test for this model
with its admissible denominator domain and allowed constraints, not a
new general theorem about Gr\"obner bases. The revised manuscript also
explains the role of the auxiliary variable $t$ and attributes the
fiber-dimension step in Theorem 4 to algebraic geometry.

\section*{Response to Reviewer 5}

\textbf{Comment.}
\begin{quote}\itshape\emergencystretch=2em
This paper presents an extension of well-known results on identifiability and informativity to a more general networked model. This reviewer has several comments:
\end{quote}

\textbf{Response.} We thank the reviewer for the detailed assessment of the model, results, and examples.

\subsection*{R5.1: Motivation for $P,Q,R,S$}
\textbf{Comment.}
\begin{quote}\itshape\emergencystretch=2em
1.- The motivation of this new model is unclear. The authors propose a model of the form $P(q)w_k=Q(q)u_k$  $Z_k=R(q)w_k+S(q)u_k$. It is not justified why there are two matrices P(q) and Q(q) instead of a single matrix. In addition, it is unusual that the dynamics for the state and the output are different. Typically, the models in network identification are of the form: $w_k=G(q)w_k+Br_k+v_k$  $y_k=Cw_k$, Natural extensions could involve considering matrices of transfer functions in B and/or C. However, a comparison or justification of the proposed model is not provided.
\end{quote}

\textbf{Response.} Section II now displays the standard network
$w_k=G(q)w_k+B(q)u_k$, $z_k=C(q)w_k$ as the special case $P=I-G$,
$Q=B$, $R=C$, $S=0$. The roles of the four matrices are stated there:
$P$ describes internal connections, $Q$ how inputs enter the nodes, $R$ how
nodes are measured, and $S$ direct input-to-output paths. The model permits
unknown entries in $Q$ and $R$, direct paths through $S$, and known relations
among entries of $P,Q,R,S$. Local physical dynamics
and couplings define the node equations while actuators define a separate
input matrix, as in the
physical-network model of Kivits and Van den Hof cited in Section II.
Normalizing the node equations can yield an $I-G$ form, but the resulting
transfer functions share factors and constraints, which the constrained-entry test
retains. Section II also gives a discrete-time descriptor example with
$P(q)=qE-A$, $Q=B$, $R=C$, and $S=D$. For singular $E$, this form
retains algebraic equations and makes prescribed zero patterns and known
coefficients explicit. Identifiability of its real matrix coefficients
requires retaining their constant-coefficient structure; the example
motivates the model form without claiming a separate descriptor
identifiability result. Filtered sensors and direct input effects
motivate $R$ and $S$; Example 3 uses repeated unknowns across matrices.

\subsection*{R5.2--3: Novelty of the main results and the path/rank link}
\textbf{Comment.}
\begin{quote}\itshape\emergencystretch=2em
2.- The main results appear to be just straightforward extensions of well-known results in informativity and identifiability of linear networks. That is why there are no mathematical proofs, except for Theorem 4, which is a bit trivial. 3.- Theorem 2 related to the link of vertex-disjoint paths with the rank of associated matrices is already well-known in identifiability of linear networks from the reference [3] of the manuscript.
\end{quote}

\textbf{Response.} We agree that the disjoint-path principle is not new;
the submitted reference [3] studies network identifiability with partial
node measurements, and [17] gives the structured-system rank result.
The revised abstract states the model class, and three bullets at the end of
Section I identify the novel model-specific contributions relative to the
cited work. Section IV cites van der Woude's structured-system result and
the submitted reference [3] as related network path literature. It explains
that the paths characterize the rank of the full closed-loop transfer function
matrix in the $P,Q,R,S$ model, including direct paths through $S$.
Theorem 2 records the direct-feedthrough extension for the present
$RP^{-1}Q+S$ form, allowing prescribed nonzero diagonal entries of $P$.
The following remark defines fixed selector matrices $Q,R$ and relates
their path-rank condition with $P=I-G$, $S=0$ to the submitted reference [3].
It also permits independent free entries in $S$. Theorem 3 handles more
general known and constrained entries through polynomial ideals. Section V then
applies a dimension test to the closed-loop map with unknown input and
measurement transfer functions. We have confined the novelty claim to this model
class and its combined rank and fiber analysis.
Example 1 uses a direct determinant calculation because its fixed
off-diagonal entries of $P$ fall outside the graph theorem's assumptions.

\subsection*{R5.4: Separation of informativity and identifiability}
\textbf{Comment.}
\begin{quote}\itshape\emergencystretch=2em
4.- Informativity and identifiability are usually two different concepts as it has been stated in the reference [13] of the manuscript. Analyzing both subjects simultaneously does not seem natural and I think they deserve separate treatment as is usually done.
\end{quote}

\textbf{Response.} The introduction now describes the two questions
separately. Sections IV and V treat them in order: informativity concerns
whether measured signals distinguish candidate predictors; network
identifiability concerns which open-loop transfer functions are determined
by the closed-loop transfer function $G_c$. The latter step can have a finite ambiguity
even when the informativity condition holds, as Example 3 shows.

\subsection*{R5.5: Length and role of the examples}
\textbf{Comment.}
\begin{quote}\itshape\emergencystretch=2em
5.- Several sections of the paper do not provide new information and give the impression that they could be easily removed. There is a full page of only examples corresponding to Section VII. In addition, Section VIII Numerical Examples does not provide interesting information. Even some of the comments are rather vague like we can compute in 11 minutes on a laptop computer .
\end{quote}

\textbf{Response.} We revised the examples around distinct points of the
analysis. Example 1 separates generic uniqueness from an exceptional
zero-edge realization; Example 2 shows predictor-input reduction from
known structure; Example 3 shows a two-point algebraic fiber and a worked
Gr\"obner basis calculation. The five-node calculations report image
dimensions and generic rank for a larger example. We replaced the vague
runtime description with the measured computation time, software
versions, processor and memory, and prepared the MATLAB supplement with
both entry scripts, three helper files, and run instructions. The
computational section now states when the path and Jacobian-rank tests may
avoid elimination and cautions that the reported runtime is not a scaling law.

\subsection*{R5.6: Subnetwork section}
\textbf{Comment.}
\begin{quote}\itshape\emergencystretch=2em
6.- The Section VI Sub-networks could be easily removed by analyzing from the beginning the case of sub-networks where the full network case can be easily encompassed.
\end{quote}

\textbf{Response.} We retained a shorter separate section because the
boundary signals measured for a subnetwork affect the predictor inputs
and the informativity calculation in the full network. Section VI writes
both direct boundary measurements and measured combinations of boundary
signals in the $P,Q,R,S$ notation. This makes clear which transfer
function is identified for each selection without introducing a second
model class.

\subsection*{R5.7: Example 3 and the Legat--Hendrickx question}
\textbf{Comment.}
\begin{quote}\itshape\emergencystretch=2em
7.- The authors seem to mix notions of identifiability, generic identifiability, local identifiability, global identifiability, etc. That is why in Example 3, they seem to contradict the conjecture in the reference [19] of the manuscript. The authors should verify this because according to the authors in [19], the conjecture has already been proven up to five unknown transfer functions in [19]. Also, in Example 3, the authors claim that there is always another pair of matrices that satisfy the identifiability. Does this imply just that the network is unidentifiable?
\end{quote}

\textbf{Response.} We no longer present Example 3 as a counterexample to
the question in submitted reference [19]. That paper considers fixed
node-selection matrices $B,C$ and independent unknown entries of $G$.
Our Example 3 has unknown entries in $Q_r$ and $R_y$, with the same $A$
appearing in both. In our model its generic fiber has exactly two
points: $A$ and $A'=-HG-A$, with $G$ and $H$ determined by $G_c$.
For proper stable $G,H,A$, the alternative $A'$ is proper and stable.
Thus the example is generically locally identifiable in the algebraic
sense but not generically uniquely network identifiable within our
broader model class. It does not resolve the question for the narrower
Legat--Hendrickx class.

\subsection*{R5.8: Visible dependence of $F$ on $X$}
\textbf{Comment.}
\begin{quote}\itshape\emergencystretch=2em
8.- The authors introduce the function F of the form $F(X)=R_yP^{-1}Q_r+S_{yr}$. This should be a function of $X$, but it does not appear in the function. What is $X$?
\end{quote}

\textbf{Response.} Section V defines $X$ as the vector of free entries of
$(P,Q_r,R_y,S_{yr})$ and writes
\[
F(X)=R_y(X)P(X)^{-1}Q_r(X)+S_{yr}(X).
\]
The subsequent discussion keeps $\lambda$ symbolic so equality means
equality of rational transfer functions.

\subsection*{R5.9: Zero edge in Example 1}
\textbf{Comment.}
\begin{quote}\itshape\emergencystretch=2em
9.- In Example 1, the authors mention that they do not have global network identifiability because there is one solution only when $G_c^1
\neq 0$. By definition, this case $G_c^1 \neq 0$ is excluded since it implies that there is no edge or connection.
\end{quote}

\textbf{Response.} Section II now distinguishes the specified structural
zero pattern from the value of a free entry at a particular realization.
A structurally free entry may vanish exceptionally. Example 1 explicitly
states that $G^1=0$ is permitted but nongeneric, and that uniqueness
fails there while generic network identifiability holds.

\subsection*{R5.10: Algebraic terminology}
\textbf{Comment.}
\begin{quote}\itshape\emergencystretch=2em
10.- The authors use concepts like affine varieties, ideals, fibers, Grobner basis, etc. that are not standard in the control systems community and for this reason, these concepts should be clearly introduced in the manuscript to improve readability.
\end{quote}

\textbf{Response.} The revised text explains the ideal generated by
polynomials, polynomial zero sets (varieties), irreducibility,
dimension, the Zariski closure, fibers, and elimination where each is
used. Section IV explains how a lexicographic Gr\"obner basis isolates
the rank-deficient set. Section V and Example 3 show the corresponding
image and fiber calculations.

\subsection*{R5.11: Typos and references}
\textbf{Comment.}
\begin{quote}\itshape\emergencystretch=2em
11.- There are many typos in the paper like the publication year (2913) in the reference [21].
\end{quote}

\textbf{Response.} We corrected notation, grammar, and terminology
throughout the manuscript. The cited bibliography year is 2013.
The revised version uses $z_k$ for the measured output signal,
$q$ for the shift operator, and $\lambda$ for a complex argument.
We also corrected Dankers's name in the bibliography and consolidated
two citations to different editions of the same algebraic-geometry book.

\section*{Materials accompanying the revision}

The source manuscript is \path{main.tex}, and the prepared MATLAB archive
is \path{supplementary_matlab_code.zip}. The archive contains the two
five-node example scripts, their three helper functions, and a README
specifying the commands and software dependencies. The manuscript points
to this archive in Section VIII.

\end{document}
